% Section 2: Software description
\section{Software description}
\label{sec:software-description}

FHOPS is organized around the core operational scheduling question: given a set of blocks, machine resources, shift calendars, and mobilization constraints, what assignments should be made, in which sequence, and with what expected production and cost consequences. The software architecture is therefore defined by planning entities and decision logic first, and by implementation boundaries second.

\subsection{Operational entities and planning context}
The platform uses a typed scenario contract allowing solver results to be compared on fixed operational inputs:
\begin{itemize}
  \item \textbf{Blocks and windows}: harvest volumes, treatment windows, and system requirements.
  \item \textbf{Machines and roles}: role-specific production rates, availability, and compatibility.
  \item \textbf{Shift calendar}: day-shift time grid used for assignment and production accounting.
  \item \textbf{Landings and mobilization}: movement penalties and daily landing capacities.
  \item \textbf{Constraint metadata}: sequencing prerequisites, role activation conditions, and completion balances.
\end{itemize}

FHOPS ships a public BC ground-based reference ladder (tiny7, small21, med42) that scales from 7/2 to 21/6 to 42/12 (days/blocks) while keeping one operational archetype.

This addresses a central review finding \cite{jaffray2025systematic}: transfer fails when assumptions are embedded in bespoke code rather than explicit scenario data. Generalizability is supported structurally through this scenario contract, which separates operational assumptions from solver logic. Empirical validation across additional system families remains future work.

\subsection{Solution methods for operational decisions}
FHOPS provides complementary solver families for different planning horizons and time budgets:
\begin{itemize}
  \item \textbf{Metaheuristic solvers (SA/ILS/Tabu)} provide practical schedule candidates for larger instances where rapid re-solving is needed.
  \item \textbf{Mixed-integer linear programming (MILP)} provides an exact baseline where tractable, supporting verification and sensitivity analysis.
  \item \textbf{Playback evaluation} re-simulates solved schedules under deterministic and stochastic conditions to quantify utilization, continuity, and disruption effects.
\end{itemize}

This reflects established forestry operations research (OR) practice: exact models support formulation clarity and benchmark bounds, while heuristics support larger-scale operational runs \cite{paradis2018bilevel,nelson2003forestmodels}.

Warm-start support is implemented for MIP, but this paper does not claim decisive wall-clock gains on med42/large84 stress contexts.

\subsection{Operational MILP and equation-to-code traceability}
The deterministic operational MILP is presented in canonical form below. Equation blocks are labeled as OBJ, E1--E11, and D1.

% AUTO-GENERATED from fhops_operational_formulation.md -- do not edit directly.
FHOPS' deterministic operational solver is defined on a day-shift grid.
It maximizes delivered production and penalizes unmet volume,
landing-capacity surplus, and machine transitions. Equation blocks below
are the canonical manuscript/docs reference and map directly to
implementation.

\textbf{Problem statement.} Given blocks, machine roles, calendars,
windows, and landing capacities, choose machine-block assignments and
per-shift production to maximize weighted output under feasibility and
sequencing constraints.

\textbf{Sets and indices.}

\begin{itemize}
\tightlist
\item
  \(m \in \mathcal{M}\): machines.
\item
  \(b \in \mathcal{B}\): blocks.
\item
  \(s=(d,\sigma) \in \mathcal{S}\): shift slots (day \(d\), shift label
  \(\sigma\)).
\item
  \(\mathcal{R}_b\): ordered roles required on block \(b\).
\item
  \(\mathcal{L}\): landings.
\item
  \(\mathcal{P}^{\text{inv}}\): role-block pairs with upstream
  prerequisites.
\item
  \(\mathcal{P}^{\text{act}} \subseteq \mathcal{P}^{\text{inv}}\):
  role-block pairs requiring head-start activation.
\item
  \(\mathcal{P}^{\text{load}}\): loader role-block pairs.
\end{itemize}

\textbf{Parameters.}

\begin{itemize}
\tightlist
\item
  \(\bar{p}_{mb}\): production rate for machine \(m\) on block \(b\).
\item
  \(W_b\): required total block production volume.
\item
  \(A_{m,s} \in \{0,1\}\): machine availability for shift \(s\).
\item
  \(\mathbf{1}^{\text{window}}_{b,d} \in \{0,1\}\): block window
  indicator.
\item
  \(\omega^{\text{prod}},\omega^{\text{mob}},\omega^{\text{trans}},\omega^{\text{land}}\):
  objective weights.
\item
  \(\delta_{m,b',b}\): mobilization cost when machine \(m\) transitions
  from block \(b'\) to block \(b\).
\item
  \(C_{\ell}\): daily assignment capacity for landing \(\ell\).
\item
  \(\ell(b)\): landing associated with block \(b\).
\item
  \(\mathcal{U}_{r,b}\): upstream roles that must feed role \(r\) on
  block \(b\).
\item
  \(B_{r,b}\): required staged buffer volume before role \(r\) may
  activate on block \(b\).
\item
  \(Q_{r,b}\): role production capacity upper bound per shift (used for
  activation linearization).
\item
  \(q^{\text{batch}}_{r,b}\): loader batch size for loader role \(r\) on
  block \(b\).
\item
  \(\mathcal{T}_b \subseteq \mathcal{R}_b\): terminal roles for block
  \(b\).
\end{itemize}

\textbf{Decision variables.}

\begin{itemize}
\tightlist
\item
  \(x_{m,b,s} \in \{0,1\}\): 1 if machine \(m\) is assigned to block
  \(b\) in shift \(s\).
\item
  \(p_{m,b,s} \ge 0\): production by machine \(m\) on block \(b\) in
  shift \(s\).
\item
  \(z_{r,b,s} \ge 0\): aggregated role-level production for role \(r\)
  on block \(b\) in shift \(s\).
\item
  \(y_{m,b',b,s} \in \{0,1\}\): transition indicator from previous-shift
  block \(b'\) to current block \(b\) (non-initial shifts).
\item
  \(I^{\text{start}}_{r,b,s} \ge 0\): staged inventory available at
  start of shift \(s\) for role \(r\) on block \(b\).
\item
  \(I_{r,b,s} \ge 0\): staged inventory at end of shift \(s\) for role
  \(r\) on block \(b\).
\item
  \(g_{r,b,s} \in \{0,1\}\): role activation indicator for buffered
  downstream roles.
\item
  \(n_{r,b,s} \in \mathbb{Z}_{\ge 0}\): loader batch count for loader
  role-block pair \((r,b)\).
\item
  \(u_{r,b,s} \ge 0\): loader partial remainder volume.
\item
  \(L_b \ge 0\): leftover unmet block volume slack.
\item
  \(S_{\ell,d} \ge 0\): landing daily surplus slack.
\end{itemize}

\textbf{Objective.}

FHOPS maximizes weighted terminal production and subtracts penalty
terms:

\[
\begin{aligned}
\max\; &\omega^{\text{prod}}\!\sum_{b\in\mathcal{B}}\sum_{r\in\mathcal{T}_b}\sum_{s\in\mathcal{S}} z_{r,b,s}
- \omega^{\text{prod}}\!\sum_{b\in\mathcal{B}} L_b \\
&- \omega^{\text{land}}\!\sum_{\ell\in\mathcal{L}}\sum_{d\in\mathcal{D}} S_{\ell,d} \\
&- \omega^{\text{mob}}\!\sum_{m,b',b,s} \delta_{m,b',b}\, y_{m,b',b,s}
- \omega^{\text{trans}}\!\sum_{m,b',b,s} y_{m,b',b,s}.
\end{aligned}
\]

If a block has no terminal-role metadata, implementation falls back to
machine-level production sums for the production-reward term.

\textbf{Constraints.} For traceability, the objective is tagged as
\textbf{OBJ}, constraints as \textbf{E1--E11}, and domain declarations
as \textbf{D1}.

Machine assignment feasibility (\textbf{E1}):

\[
\sum_{b\in\mathcal{B}} x_{m,b,s} \le A_{m,s}
\qquad \forall m\in\mathcal{M},\; s\in\mathcal{S}.
\]

Role compatibility (\textbf{E2}, machines can only work roles allowed by
the block's assigned harvest system):

\[
x_{m,b,s}=0 \quad \text{if role}(m)\notin\mathcal{R}_b.
\]

Production upper bound per assignment (\textbf{E3}):

\[
p_{m,b,s} \le \bar{p}_{mb}\,x_{m,b,s}
\qquad \forall m,b,s.
\]

Block window enforcement (\textbf{E4}):

\[
x_{m,b,s}=0 \quad \text{if } \mathbf{1}^{\text{window}}_{b,d}=0 \text{ for } s=(d,\sigma).
\]

Role-production aggregation (\textbf{E5}):

\[
z_{r,b,s} = \sum_{m\in\mathcal{M}(r)} p_{m,b,s}
\qquad \forall (r,b), s.
\]

Transition linking (\textbf{E6}, for non-initial shifts only):

\[
y_{m,b',b,s} \le x_{m,b',\operatorname{prev}(s)},
\qquad
y_{m,b',b,s} \le x_{m,b,s},
\] \[
y_{m,b',b,s} \ge x_{m,b',\operatorname{prev}(s)} + x_{m,b,s} - 1.
\]

Role inventory start and balance for prerequisite-driven downstream
roles (\textbf{E7}):

\[
I^{\text{start}}_{r,b,s}=
\begin{cases}
0, & s \text{ is first shift}\\
I_{r,b,\operatorname{prev}(s)}, & \text{otherwise}
\end{cases}
\qquad \forall (r,b)\in\mathcal{P}^{\text{inv}}, s,
\]

\[
I_{r,b,s}=I^{\text{start}}_{r,b,s}+\sum_{u\in\mathcal{U}_{r,b}} z_{u,b,s}-z_{r,b,s}
\qquad \forall (r,b)\in\mathcal{P}^{\text{inv}}, s,
\]

\[
z_{r,b,s} \le I^{\text{start}}_{r,b,s}
\qquad \forall (r,b)\in\mathcal{P}^{\text{inv}}, s.
\]

Head-start activation for buffered downstream roles (\textbf{E8}):

\[
z_{r,b,s} \le Q_{r,b}\,g_{r,b,s}
\qquad \forall (r,b)\in\mathcal{P}^{\text{act}}, s,
\]

\[
I_{r,b,\operatorname{prev}(s)} \ge B_{r,b}\,g_{r,b,s}
\qquad \forall (r,b)\in\mathcal{P}^{\text{act}}, s,
\]

\[
\sum_{m\in\mathcal{M}(r)} x_{m,b,s} \le |\mathcal{M}(r)|\, g_{r,b,s},
\qquad
g_{r,b,s} \le \sum_{m\in\mathcal{M}(r)} x_{m,b,s}
\qquad \forall (r,b)\in\mathcal{P}^{\text{act}}, s.
\]

Loader batching (\textbf{E9}):

\[
z_{r,b,s}=q^{\text{batch}}_{r,b}\,n_{r,b,s}+u_{r,b,s}
\qquad \forall (r,b)\in\mathcal{P}^{\text{load}}, s,
\]

\[
0 \le u_{r,b,s} \le q^{\text{batch}}_{r,b}
\qquad \forall (r,b)\in\mathcal{P}^{\text{load}}, s.
\]

Block completion balance with leftover slack (\textbf{E10}):

\[
\sum_{r\in\mathcal{T}_b}\sum_{s\in\mathcal{S}} z_{r,b,s} + L_b = W_b
\qquad \forall b\in\mathcal{B}.
\]

Landing daily assignment capacity with surplus slack (\textbf{E11}):

\[
\sum_{b:\,\ell(b)=\ell}\sum_{m\in\mathcal{M}}\sum_{\sigma:(d,\sigma)\in\mathcal{S}} x_{m,b,(d,\sigma)}
\le C_{\ell} + S_{\ell,d}
\qquad \forall \ell\in\mathcal{L},\; d\in\mathcal{D}.
\]

Domain restrictions (\textbf{D1}):

\[
x, y, g \in \{0,1\},\quad n \in \mathbb{Z}_{\ge 0},\quad p,z,I^{\text{start}},I,u,L,S \ge 0.
\]

\textbf{Implementation mapping (equation blocks to code).}

Primary modules: \texttt{src/fhops/model/milp/operational.py},
\texttt{src/fhops/model/milp/data.py}. Primary tests:
\texttt{tests/model/test\_operational\_milp.py},
\texttt{tests/model/test\_operational\_driver.py}.

Equation-block mapping:

\begin{itemize}
\tightlist
\item
  \textbf{OBJ}: \texttt{model.objective} with \texttt{prod\_weight},
  \texttt{landing\_weight}, \texttt{mobilisation\_weight},
  \texttt{transition\_weight}
\item
  \textbf{E1}: \texttt{model.machine\_capacity}
  (\texttt{machine\_capacity\_rule})
\item
  \textbf{E2}: \texttt{model.role\_compatibility}
  (\texttt{role\_compatibility\_rule})
\item
  \textbf{E3}: \texttt{model.production\_cap} (\texttt{prod\_cap\_rule})
\item
  \textbf{E4}: \texttt{model.block\_windows} (\texttt{window\_rule})
\item
  \textbf{E5}: \texttt{model.role\_prod\_balance}
  (\texttt{role\_prod\_balance\_rule})
\item
  \textbf{E6}: \texttt{model.transition\_prev},
  \texttt{model.transition\_curr}, \texttt{model.transition\_link}
\item
  \textbf{E7}: \texttt{model.inventory\_start\_eq},
  \texttt{model.inventory\_balance}, \texttt{model.inventory\_guard}
\item
  \textbf{E8}: \texttt{model.activation\_prod},
  \texttt{model.head\_start}, \texttt{model.role\_active\_upper},
  \texttt{model.role\_active\_lower}
\item
  \textbf{E9}: \texttt{model.loader\_batch},
  \texttt{model.loader\_partial\_cap}
\item
  \textbf{E10}: \texttt{model.block\_balance}
  (\texttt{block\_balance\_rule}) + \texttt{model.leftover}
\item
  \textbf{E11}: \texttt{model.landing\_capacity}
  (\texttt{landing\_capacity\_rule}) + \texttt{model.landing\_surplus}
\item
  \textbf{D1}: variable domain declarations for \texttt{model.x},
  \texttt{model.y}, \texttt{model.role\_active}, \texttt{model.loads},
  \texttt{model.prod}, \texttt{model.role\_prod},
  \texttt{model.inventory\_start}, \texttt{model.inventory},
  \texttt{model.loader\_partial}, \texttt{model.leftover}, and
  \texttt{model.landing\_surplus}
\end{itemize}

This formulation is the canonical mathematical reference for FHOPS
operational MILP documentation and companion modelling manuscripts.


Each block maps to named implementation objects and regression tests. The shared formulation source keeps the manuscript and documentation synchronized for auditability.
For practical review, the path is explicit: equation labels map to Pyomo objects in the operational builder, then to tests that check sequencing logic, bundle normalization, and driver integration. This design and structure reduces ambiguity when readers compare mathematical intent with executable behaviour and helps maintainers detect formulation drift when model code evolves.

\subsection{Architecture and implementation boundaries}
Internally, FHOPS separates four concerns:
\begin{enumerate}
  \item \textbf{Scenario layer}: ingestion and validation of planning entities.
  \item \textbf{Optimization layer}: heuristic and MILP solvers sharing common scenario inputs.
  \item \textbf{Evaluation layer}: playback and key performance indicators (KPI) computation for solved schedules.
  \item \textbf{Reporting layer}: generation of tables/figures used in manuscript and documentation outputs.
\end{enumerate}

The software is available through Python modules and command-line wrappers. Interface choice is secondary to preserving equivalent, traceable evidence across runs. FHOPS is released under an MIT license, with source code, datasets, and experiment pipelines in a public repository to support reuse and independent verification.

\subsection{Reproducible evidence pipeline}
Manuscript evidence comes from versioned benchmark, tuning, playback, costing, and scaling assets produced by the same platform layers. Run logs and artifact paths provide provenance for each reported value.

This arrangement is deliberate in the manuscript sequence: the review identifies missing reproducibility infrastructure \cite{jaffray2025systematic}, this paper contributes it as reusable software, and companion modelling studies can focus on method questions without redefining data contracts or reporting conventions.
